\section{Methods}
\subsection{Study design and anatomical data}
The anatomical forward simulations used the original BrainWeb normal phantom, with its noise-free, intensity-uniform T1-weighted image and discrete tissue labels \citep{collins1998,brainweb}. The input is an anatomical substrate, not a longitudinal glioma dataset. No patient data were used to calibrate these phantom experiments; the separate PAT16 input-based pilot is described below. The native 1 mm volumes contain $181\times217\times181$ voxels. The primary simulation used a 2 mm grid of $91\times109\times91$ voxels; intensities were resampled linearly and labels by nearest neighbor. Arrays use canonical RAS coordinates. Gray matter, white matter and the phantom's glial class form the transport domain $\Omega$. CSF and other tissues are impermeable in every matched comparison. Data preparation and checksums are in Appendix~\ref{app:data}. Synthetic inference was evaluated first in a homogeneous phantom, then on the full BrainWeb anatomy and finally on PAT16 anatomy with directional transport. The phantom recovery experiments use dilute known seeds; PAT16 uses a known lesion-based initial density. All observations fitted for recovery are synthetic, and none calibrates the dense-seed coupled comparison or patient growth rates.

\subsection{Reference model: Fisher--KPP}
Let $c(\bm{x},t)$ denote viable tumor-cell density in cells/mm$^3$. The reference, M0, is
\begin{equation}
 \partial_t c=\nabla\!\cdot\!\left(D(\bm{x})\nabla c\right)+\rho(\bm{x})c\left(1-\frac{c}{K}\right),\qquad\bm{x}\in\Omega .\label{eq:kpp}
\end{equation}
$D$ is effective cell migration in mm$^2$/day, $\rho$ is growth in day$^{-1}$, and $K$ is a density scale. $D$ is scalar with separate white- and gray-matter values in the BrainWeb experiments; the separate PAT16 pilot replaces it by a trace-matched tensor. Cell and, when present, oxygen transport have homogeneous no-flux boundaries. The matched reference uses the same seed, coefficients, anatomy and output times as the extension.

\subsection{Oxygen-limited and three-field extensions}
The coupled model, M2, adds relative oxygen availability $o(\bm{x},t)$ and immobile necrotic material $z(\bm{x},t)$ in tumor-cell equivalents/mm$^3$. Define
\begin{equation}
 \phi=\frac{c+z}{K},\quad g(o)=\min\left(1,\frac{o}{\theta_p}\right),\quad h(o)=\pos{1-\frac{o}{\theta_d}},\quad p(\phi)=\frac{1}{1+\alpha\phi},\label{eq:closures}
\end{equation}
where $[a]_+=\max(a,0)$ and $0<\theta_d<\theta_p<1$. The system is
\begin{align}
 \partial_t c&=\nabla\!\cdot(D\nabla c)+\rho g(o)c\pos{1-\phi}-\delta h(o)c,\label{eq:viable}\\
 \partial_t z&=\delta h(o)c-\lambda_z z,\label{eq:dead}\\
 0&=\ell_o^2\nabla^2o+p(\phi)(1-o)-\beta(c/K)o.\label{eq:oxygen}
\end{align}
Oxygen is solved at initialization and after every cell step. Its quasi-steady treatment imposes faster oxygen equilibration than tumor evolution, an assumption not tested by these data. Normal-tissue metabolism is implicit in the baseline $o=1$. Distributed restoration replaces explicit vessels, and $p$ is a phenomenological supply coefficient rather than measured blood flow. Viable cells consume oxygen. Necrotic material affects supply and available growth capacity but consumes no oxygen.

The two-field model M1 sets $\delta=0$ and starts with $z=0$, so oxygen only modulates growth. M2 includes death transfer and optional clearance. All examples start with $z=0$. With $\beta=0$, oxygen is one and M2 reduces to M0 for uniform $K$ and the stated initial condition. Setting $\alpha=0$ holds restoration fixed while retaining uptake. These nested cases provide implementation checks and ablations.

$K$ regulates births rather than imposing hard packing on $c+z$. Diffusion can enter an already occupied location. Consequently, $z/K$ and $\phi$ can exceed one: $z/K$ is not a physical necrotic volume fraction. No healthy-cell population, edema, angiogenic state, tissue displacement, directed hypoxic migration or treatment response is included.

\subsection{Initialization, parameters and numerical experiments}
The initial density was a tissue-restricted Gaussian centered at RAS $(-30,25,30)$ mm, with standard deviations of 6 mm in each direction, a peak of $80{,}000$ cells/mm$^3$, and truncation at three standard deviations. This dense seed was selected to produce an oxygen-limited interior early in the demonstration. Simulations cover 90 days, saving every 3 days. All biological values in Table~\ref{tab:parameters} were assigned working values, not fitted estimates or established patient-level priors.

\begin{table}[tbp]
\centering\small
\caption{Assigned parameters for the primary comparison. All are illustrative. Numerical tolerances and initialization details are in the appendices.}\label{tab:parameters}
\begin{tabularx}{\textwidth}{l Y l l}
\toprule Symbol & Interpretation & Value & Unit\\\midrule
$D_{\rm WM}$ & White-matter migration & 0.5 & mm$^2$/day\\
$D_{\rm GM}$ & Gray/glial-matter migration & 0.1 & mm$^2$/day\\
$\rho$ & Growth coefficient & 0.04 & day$^{-1}$\\
$K$ & Density scale & $100{,}000$ & cells/mm$^3$\\
$\ell_o$ & Effective oxygen exchange length & 0.2 & mm\\
$\beta$ & Relative tumor oxygen uptake & 4 & dimensionless\\
$\alpha$ & Supply suppression by total burden & 4 & dimensionless\\
$\theta_p$ & Oxygen for unrestricted proliferation & 0.3 & relative oxygen\\
$\theta_d$ & Hypoxic-death threshold & 0.1 & relative oxygen\\
$\delta$ & Maximum death-transfer rate & 0.5 & day$^{-1}$\\
$\lambda_z$ & Necrotic-material clearance & 0 & day$^{-1}$\\\bottomrule
\end{tabularx}
\end{table}

The primary M0--M2 comparison used $\Delta x=2$ mm and maximum $\Delta t=0.25$ day. Additional 2 mm runs evaluated M1, M2 with $\alpha=0$, and M2 with $\beta=0$. Temporal assessment repeated M2 with $\Delta t=0.125$ and 0.0625 day. Native-grid assessment repeated M2 and matched M0 at 1 mm with a 0.125-day step. These nine model/step/grid cases are controlled simulations, not biological replicates. The grid comparison changes discretized anatomy as well as the numerical mesh, and is not a pure spatial-order experiment.

Cell transport uses conservative face fluxes with harmonic interface coefficients. Explicit diffusion is followed by exact frozen-coefficient logistic birth, hypoxic-death transfer and clearance substeps. Oxygen uses a compact no-flux stencil and preconditioned conjugate gradients. The method is first order in time. A diffusion positivity bound controls substeps but does not guarantee reaction accuracy. Details are in Appendix~\ref{app:numerics}.

\subsection{Outcome measures and interactive visualization}
We recorded viable count $N_c=\int_\Omega c\,\dd V$, necrotic equivalents $N_z=\int_\Omega z\,\dd V$, peak density, minimum tissue oxygen, and volumes above $c/K=0.2$ and 0.8. Necrotic threshold volume uses $z/K\geq0.2$; low-oxygen volume uses $o<\theta_d$. These are model-defined quantities, not validated MRI detection thresholds. With closed boundaries,
\begin{equation}
 \frac{\dd}{\dd t}(N_c+N_z)=\int_\Omega\rho g(o)c\pos{1-\phi}\,\dd V-\lambda_zN_z .\label{eq:balance}
\end{equation}
Death transfer cancels. With zero clearance and matched initial states, the M0 viable-count minus M2 live-plus-dead-count difference measures the cumulative birth deficit within the numerical balance.

The Jupyter interface exposes biological and numerical controls and retains Fisher--KPP. Four fields share time and slice coordinates, with axial, coronal and sagittal views. Live 3D histories support arbitrary slicing; compact saved results retain seed-plane histories and final volumes. MP4/GIF exports share a frame schedule. A slower movie spans 20 seconds with a 6-second day-90 hold. Display settings do not alter fields. The paper uses lossless arrays and displays the full necrotic range.

\subsection{Homogeneous recovery of diffusion and proliferation}\label{sec:recovery_methods}
To test two-parameter estimation independently of the anatomical comparisons, we solved normalized M0 dynamics, $u=c/K$, in a homogeneous 60 mm cubic no-flux domain. Ground-truth rates were $D=\RecoveryDTrue$ mm$^2$/day and $\rho=\RecoveryRhoTrue$ day$^{-1}$. The known initial Gaussian had peak 0.2, standard deviations $(5,4,5)$ mm and center $(-3,2,0)$ mm in box-centered coordinates. It was untruncated, with analytically integrated cell averages; this differs from the dense tissue-restricted BrainWeb seed. The synthetic signal applied a Gaussian point-spread function of 2 mm standard deviation and averaged disjoint 6 mm cubic voxels. Independent zero-mean Gaussian noise, with known standard deviation 0.03 in normalized signal units, was added after imaging. Measurements were not clipped. This is a calibrated cellularity proxy, not a raw MRI signal or a claim that water ADC measures cell migration.

The generator used 1 mm cells and a maximum 0.0625-day step; fitting used 2 mm cells and a 0.25-day step with the same conservative diffusion/logistic kernel in float64. Different discretizations reduce numerical self-consistency bias, while the shared equations and observation model retain favorable structural assumptions. We measured their observation-space discrepancy explicitly. Scans at days 20, 40 and 60 supplied 3,000 scalar training measurements; days 75 and 90 were withheld from fitting and uncertainty estimation. The normalized initial state, age origin, signal calibration, blur and noise were fixed and known. Uniform $K$ cancels from normalized M0 dynamics; these recovery experiments do not identify absolute capacity or cell counts.

Bounded Gaussian maximum likelihood estimated $(\log D,\log\rho)$ using nonlinear least squares \citep{virtanen2020}, with three starting guesses and no probability prior. Local sensitivity-based marginal 95\% confidence intervals were compared with percentile intervals from \RecoveryBootstraps{} parametric bootstrap fits \citep{efron1979}. Bootstrap parameter pairs generated pointwise confidence bands for the noise-free mean signal; adding independent measurement noise generated pointwise new-scan prediction intervals. A separate \RecoveryTrials{}-dataset noise experiment, generated at the known rates on the fine grid, assessed recovery bias, RMSE and local-interval coverage with exact binomial uncertainty. This audits local intervals, not bootstrap coverage. A time-zero-only negative control tested rejection of unidentifiable rates. Equations, bounds, diagnostics and resampling details are in Appendix~\ref{app:recovery}.

\subsection{Three-rate recovery on BrainWeb and anatomy sensitivity}\label{sec:bw_methods}
The second inference experiment estimated $D_{\rm WM}$, $D_{\rm GM}$ and a common $\rho$ in normalized M0 dynamics on the full BrainWeb parenchyma. Ground truth was $(0.500,0.100,0.040)$, in mm$^2$/day, mm$^2$/day and day$^{-1}$. Label 3 used $D_{\rm WM}$ and labels 2/8 used $D_{\rm GM}$; CSF and extracerebral tissue were impermeable. No ordering or fixed ratio linked the diffusion rates. The known seed was an untruncated, point-sampled, tissue-restricted Gaussian of peak 0.2 and 6 mm widths at RAS $(-30,25,30)$ mm. The T1 image supplied anatomical context rather than measured tumor signals.

The primary generator and inverse model shared the full 3 mm tissue grid ($61\times73\times61$, 58,669 active voxels), with maximum steps 0.125 and 0.5 day, respectively. This isolates temporal mismatch but retains a shared spatial discretization; it is a favorable conditional recovery test. A fixed 60 mm observation cube centered at $(-30,24,30)$ mm supplied $10^3$ voxels per scan. The PDE domain was not cropped. A known 2 mm Gaussian blur preceded physical-coordinate interpolation and approximate 6 mm voxel averaging. Independent Gaussian noise of known standard deviation 0.02 was added after imaging. All 3,000 measurements at days 20/40/60 were fitted; days 75/90 were withheld. Different noise levels and seeds prevent interpreting comparisons with the homogeneous experiment as an isolated effect of anatomy.

The three-log-rate likelihood, local covariance and percentile bootstrap followed the preceding procedure, with \BWBootstraps{} bootstrap fits and \BWTrials{} independent fine-time-step noise realizations. The latter audited local intervals, not bootstrap coverage. Three starts and sensitivity singular values checked local estimation. Forecasts were assessed over the full observation field of view (FOV) and a prespecified 18 mm radius around the known seed, containing 120 observation-voxel centers. This region was fixed before generating signal or noise.

For a separate sensitivity test, the generator used the full 2 mm BrainWeb grid with the same physical observation geometry, analytic seed and 0.125-day step. The 3 mm inverse model fitted both its noise-free signal and a noisy signal using the primary noise draw for a paired comparison. The test changes tissue geometry, mesh, seed sampling and discrete imaging; it does not isolate spatial order or estimate a distribution of segmentation errors. No confidence interval was assigned to the deterministic noise-free rate shift. Appendix~\ref{app:bw_recovery} supplies implementation and resampling details.

\subsection{PAT16: MRI-informed directional transport and static connectivity}\label{sec:pat16_methods}
A separate pilot reused preoperative sub-PAT16 derivatives from the BTC dataset and linked digital-twin workflow \citep{aerts2022,lundervold2026twin}. It used processed T1 anatomy, named CSF/gray-/white-matter probability maps, a lesion mask, a registered diffusion tensor, a cortical/subcortical atlas and resting-state BOLD with confounds. Existing fMRIPrep and MRtrix-derived inputs were read in place \citep{esteban2019,tournier2019}; preprocessing was not rerun. T1-Gd, T2 and FLAIR were not used. This is a hypothetical untreated 90-day evolution initialized from a preoperative lesion, not an estimate of tumor age, observed postoperative progression or patient growth rates.

With $u=c/K$, both variants solve $\partial_t u=\nabla\cdot(\mathbf D\nabla u)+\rho u(1-u)$. The isotropic reference uses $\mathbf D=d\mathbf I$, with separate assumed white-/gray-matter values. In eligible white matter, a unit principal dMRI axis $\bm v$ defines
\begin{equation}
 \mathbf D=D_\perp\mathbf I+(D_\parallel-D_\perp)\bm v\bm v^{\mathsf T},\qquad
 D_\perp=\frac{3d}{r+2},\quad D_\parallel=rD_\perp .\label{eq:pat16_tensor}
\end{equation}
Thus $\operatorname{tr}(\mathbf D)/3=d$ in both variants: measured water diffusion informs orientation, not the magnitude of cell migration. Positive tensor eigenvalues and $0.2<\mathrm{FA}\leq1$ define eligible white matter; other tissue remains isotropic. Tensor components are resampled and rigidly reoriented before eigendecomposition. The full 3 mm tissue domain uses conservative no-flux tensor transport, retaining mixed derivatives (Appendix~\ref{app:pat16}). Oxygen, necrosis, displacement and treatment are absent from this pilot.

We generated \PATScenarios{} paired Latin-hypercube scenarios (seed 16026), varying white-/gray-matter migration, proliferation, anisotropy, initial amplitude, lesion threshold and a small global directional perturbation. All distributions are explicit assumptions (Table~\ref{tab:pat16_priors}); they were not learned from PAT16 or transferred from synthetic recovery as patient estimates. Each pair shares its initial state and all rate draws. States were saved every 15 days through day 90. Outcomes are $\int u\,\dd V$ in mL of equivalent normalized burden, volume with $u\geq0.16$, and mean simulated density in atlas regions. None is a calibrated MRI tumor volume, absolute cell count or functional-risk measure. Empirical 5th/50th/95th percentiles describe prior/sensitivity propagation, not confidence or posterior credible intervals.

Static functional connectivity was recomputed from actual preprocessed BOLD after atlas resampling, censoring and joint nuisance regression. The 82 regions comprise 68 Desikan--Killiany cortical and 14 subcortical regions \citep{desikan2006}. Slice overlays project at most 35 strong correlations with $|r|\geq0.35$ whose ROI centers lie within a $\pm12$ mm slab. FC is not evolved or used to alter proliferation; it provides spatial context, not axonal paths or predicted functional loss. The executed notebook synchronizes time, plane, slice and density-quantile controls across isotropic, anisotropic, paired-difference and threshold-exceedance views.

\subsection{Four-parameter synthetic recovery on PAT16 anatomy}\label{sec:pat16_recovery_methods}
A third inference experiment fitted synthetic longitudinal signals on the same real preoperative PAT16 anatomy and dMRI directions. It jointly estimated separate $D_{\rm WM}$ and $D_{\rm GM}$, shared $\rho$ and one migration ratio $r=D_\parallel/D_\perp$ shared across eligible white matter in Equation~\eqref{eq:pat16_tensor}. Known generating values were $(0.500,0.100,0.040,2.500)$, with diffusion in mm$^2$/day, proliferation in day$^{-1}$ and dimensionless $r$. Mean diffusivity and anisotropy were separate parameters; no fixed tissue ratio, ordering between WM and GM diffusivities or probability prior was imposed. Gray matter and non-eligible white matter remained isotropic. Tissue geometry, directions, eligibility and normalized initialization were fixed. Uniform $K$ cancels from these normalized dynamics; neither absolute capacity nor cell counts were identified. Unlike the pilot's assumed-parameter sensitivity bands, the intervals here were fitted conditional confidence intervals. Neither experiment estimated rates from measured patient tumor signals.

The full 3 mm tissue domain was evolved, with double-precision snapshots for likelihood differentiation. The known initial density used lesion threshold 0.5, smoothing with standard deviation 0.7 grid voxels and amplitude 0.6. A prespecified 72 mm observation cube centered on the tissue-restricted initial lesion centroid supplied $12^3=1{,}728$ voxels per scan. Gaussian blur with 2 mm standard deviation preceded 6 mm voxel averaging with $3^3$ midpoint quadrature samples and physical-coordinate interpolation. Independent Gaussian noise of known standard deviation 0.015 was added after imaging, without clipping. The synthetic cellularity scale was assumed known; these signals were not T1, FLAIR or measured patient cellularity. BOLD and functional connectivity did not enter the likelihood.

The generator and inverse shared spatial anatomy, the tensor stencil, initialization and observation model, with maximum steps of 0.25 and 0.5 day, respectively. All 5,184 scalar measurements at days 20/40/60 were fitted; days 75/90 were withheld from optimization and uncertainty estimation. Three starts minimized standardized Gaussian residuals in four log parameters. Numerical bounds were $(0.05,0.02,0.01,1.0)$ to $(1.0,0.6,0.09,3.5)$; the tensor family assumes prolate migration ($r\geq1$), and the upper bound is a monotonicity restriction of this stencil. Full-rank Jacobian information supplied local marginal 95\% log-Wald intervals, using the known noise variance without residual-variance rescaling.

We compared local intervals with percentile intervals from \PRBootstraps{} parametric bootstrap refits of noisy early scans generated at the fitted parameters. A separate \PRTrials{}-dataset audit added fresh noise at the known, finer-time generating parameters to evaluate bias, RMSE and marginal local-interval coverage with exact binomial uncertainty. Every attempted trial remained in the coverage denominator; failed or boundary local intervals counted as misses. Marginal bootstrap parameter-interval coverage and joint four-parameter coverage were not audited. Bootstrap trajectories produced pointwise signal intervals. For new-scan prediction, we compared the original 100 noise-augmented trajectories with \PPDraws{} predictive draws resampling these trajectories, and with numerical quantiles of their exact Gaussian-noise mixture. The latter supplies the final prediction intervals, removing noise-tail sampling error while retaining finite-bootstrap error. All interval levels are nominal 95\%. Forecasts were evaluated over the full cube and a prespecified 24 mm sphere around its center, containing 280 observation-voxel centers. A separate \PPTrials{}-dataset predictive audit regenerated early and withheld observations independently at the known finer-time signal. Each dataset was refitted from the default start, followed by 100 new bootstrap refits and all three predictive constructions. Thus the audit refitted uncertainty in each repetition rather than testing repeated noise against one fixed band. We report coverage, width and the proper 95\% interval score \citep{gneiting2007}. Monte Carlo intervals for mean spatial coverage use independent datasets as replicates; a fixed central voxel additionally receives exact binomial coverage limits. Withheld observations were used only for evaluation; no tail or width correction was tuned against them. Noise-free matched-model recovery, a finer-time noise-free fit and a time-zero-only negative control checked implementation and local identifiability. Appendix~\ref{app:pat16_recovery} gives the operator, optimization and reproducibility details.


\subsection{Proposed measurement-based estimation and validation protocol}
Beyond the synthetic experiments, modality $m$ at time $t_j$ would be represented by
\begin{equation}
 y_{mj}=\mathcal H_m[c_{\boldsymbol\theta}(t_j),o_{\boldsymbol\theta}(t_j),z_{\boldsymbol\theta}(t_j);\boldsymbol\eta_m]+\varepsilon_{mj},\label{eq:observation}
\end{equation}
where $\mathcal H_m$ includes contrast, resolution, partial volume and the biological-state-to-measurement link. $\boldsymbol\theta$ includes biology and initialization; $\boldsymbol\eta_m$ contains observation parameters. Calibration must incorporate noise, registration and initial-state uncertainty. An objective and identifiability restrictions are in Appendix~\ref{app:inference}.

The proposed sequence estimates a small growth--migration subset from serial MRI, introduces oxygen parameters when additional modalities support them, and calibrates cell-scale and necrosis mappings using independent samples. Numerical verification, synthetic recovery, withheld-time prediction, tissue agreement and external-cohort performance are evaluated separately. M0 and M2 must be fitted separately to identical calibration data before predictive comparison.
